\chapter{Külső szolgáltatások és egress policy}

Egy valódi webalkalmazás ritkán marad teljesen önálló. Fizetési szolgáltató, directory, belső API, térképszolgáltatás vagy más HTTPS provider külön trust boundary, mert itt maga az alkalmazás kezdeményez hálózati kapcsolatot. Az RWLang ezt explicit capabilityként kezeli, nem ambient hálózati jogosultságként.

\section{A jelenlegi RWLang-határ}

Az RWLang szándékosan nem ad korlátlan \code{http.get(userUrl)} jellegű API-t. Ehelyett a source named integrationt deklarál, a handler pedig explicit capabilityt kap:

\begin{lstlisting}[language=RWLang]
integration Billing {
    egress payments
}

action fn charge(
    ctx: ActionContext,
    amount: Int
) -> Result<Json, PageError>
    uses Billing
{
    let status = Billing.postJson("/v1/charges", amount)?;
    return Ok(json(status));
}
\end{lstlisting}

A deklaráció csak trusted \emph{nevet} tesz láthatóvá. Host, IP, port, TLS és hálózati tartomány nem lesz az alkalmazáskód authorityje. A compiler a hívást \code{net.payments} effectként rögzíti; a handler a megfelelő \code{uses Billing} capability nélkül nem indíthatja el.

A biztonságos út:

\begin{lstlisting}
RWLang business intent
        |
        v
named integration capability
        |
        v
trusted operator egress policy
        |
        v
bounded HTTPS transport
\end{lstlisting}

\section{Miért nem reprezentálható arbitrary URL?}

A user-controlled fetch target klasszikus SSRF kockázat. Loopback, link-local metadata service vagy belső infrastruktúra nem válhat normál alkalmazási céllá:

\begin{lstlisting}
http://127.0.0.1/...
http://169.254.169.254/...
http://internal-db.example/...
http://[::1]/...
\end{lstlisting}

Az RWLang source a deklarált integráción belül relatív pathot választhat, de scheme-et, hostot, IP-t vagy portot nem. A path compiler által ismert rooted relative literal; absolute, protocol-relative és dinamikus destination URL tiltott.

\section{Named target: a hálózati authority konfigurálása}

A trusted egress policy köti a nyelvi target-nevet a konkrét transport authorityhez:

\begin{lstlisting}
[[target]]
name = "payments"
hosts = ["api.payments.example"]
cidrs = ["203.0.113.0/24", "2001:db8:1200::/48"]
ports = [443]
tls_required = true
max_dns_answers = 8
max_sent_bytes = 262144
max_received_bytes = 2097152
connect_timeout_ms = 5000
total_timeout_ms = 15000
\end{lstlisting}

A typed RWLang outbound surface által használt targetet a startup policy egyértelműen a támogatott endpoint formára oldja fel. Outbound effectet használó program hiányzó vagy inkonzisztens egress policyval nem indul el.

\section{A policy több réteget ellenőriz}

Az egress policy több egy hostname allowlistnél.

\subsection*{Host, port és TLS}

A destination host és port trusted policyból jön, a TLS identity verification pedig a konfigurált hostname-et használja. Alkalmazási input nem kapcsolhatja ki a TLS-t és nem választhat másik portot.

\subsection*{DNS és CIDR}

Minden A és AAAA választ ellenőrizni kell az engedélyezett tartományok ellen. Ha egyetlen answer sérti a policyt, a feloldás fail-closed módon elutasított. Ez DNS rebinding és hibás split-horizon konfiguráció ellen is véd.

\subsection*{Peer verification connect után}

TCP kapcsolat után a tényleges peer IP újra ellenőrzött. A DNS-validáció nem maradandó proof.

\section{IPv4 és IPv6 ugyanazt az authority modellt követi}

A policy dual-stack. Az IPv4-mapped IPv6 cím normalizálódik, ezért IPv4 tiltás nem kerülhető meg IPv6 reprezentációval. Loopback, link-local, ULA, multicast és unspecified tartomány implicit trustot nem kap.

\section{Typed outbound műveletek}

Az első RWLang call surface szándékosan kicsi:

\begin{lstlisting}[language=RWLang]
let status = Billing.get("/v1/status")?;
let status = Billing.postJson("/v1/charges", amount)?;
\end{lstlisting}

Fontos korlátok:

\begin{itemize}
  \item page handler outbound GET-et végezhet, outbound POST side effectet nem;
  \item JSON POST action capability;
  \item generic outbound JSON sink nem fogad \code{Secret}, \code{Sensitive}, credential-purpose vagy upload értéket;
  \item server-side egress nem ad browser network authorityt -- a generált CSP \code{connect-src 'none'} marad, amíg nincs külön browser capability;
  \item a jelenlegi typed call HTTP status code-ot ad vissza, nem automatikusan trusted response body-t.
\end{itemize}

Ez utóbbi tudatos döntés. Külső response data csak untrusted adatként, typed validation/refinement boundaryn keresztül kerülhetne domain értékké; arbitrary upstream JSON nem válik csendben trusted adattá.

\section{Méret-, protokoll- és időkorlátok}

Az upstream lehet hibás vagy ellenséges. A transport ezért DNS answer-, request byte-, response byte- és időkorlátot alkalmaz. Az operator target limitjein felül a typed RWLang status-call surface saját hard upstream-body ceilingt is használ, és elutasítja a nem támogatott transfer/content encodingot.

A teljes request részt vesz a cumulative external-I/O budgetben is. Több külön-külön megengedett network operation így sem állhat össze egyetlen korlátlan kéréssé.

\section{Secretek külső szolgáltatáshoz}

A külső credential külön secret-store boundary mögött marad:

\begin{lstlisting}
/run/secrets/rwlang/
  payments-bearer
\end{lstlisting}

A secret név nem válhat arbitrary filesystem pathszá. Credential nem kerülhet response-ba, normál logba, security eventbe vagy audit payloadba. A typed \code{Secret<T>} és a trusted Rust secret store ugyanazt a határt erősíti két külön rétegen.

\section{Retry, idempotencia és exceptional outcome}

Side-effecting POST automatikus retry-ja szándékosan nem default. Retry-szenzitív critical operation az RWLang idempotency contractját használja, ne ad-hoc retry loopot.

Critical database worknél a transaction outcome explicit:

\begin{lstlisting}[language=RWLang]
let outcome = transaction db {
    charge(tx)?;
};

match outcome {
    Committed => { return Ok(json(true)); }
    RolledBack => { fail conflict; }
    CommitUnknown => { fail conflict; }
}
\end{lstlisting}

Idempotens critical transactionnek kötelező az outcome capture és exhaustive kezelés. A \code{CommitUnknown} nem jelentheti csendben azt, hogy ``retry POST'', mert a side effect már commitolhatott.

\section{A webhook az ellenkező irányú boundary}

Outbound integration és inbound webhook két külön probléma. A webhook route deklarált verification policyt, raw-body signature verificationt, timestamp/replay-window ellenőrzést és atomic replay claimet használ még a handler authority előtt. Webhookot ne modellezz egyszerű public JSON POST-ként.

\section{Szinkron request vagy háttérmunka?}

Jelenleg nincs first-class background-job subsystem. Rövid, bounded integration maradhat request/response ciklusban. Hosszú vagy megbízhatóan retry-zandó workflowhoz explicit job/outbox design kell; ne ad-hoc thread, subprocess vagy korlátlan request-idő.

\section{Minimum tesztelés}

Bizonyítsd legalább:

\begin{itemize}
  \item unknown integration vagy hiányzó \code{uses} capability compile-time reject;
  \item dynamic és absolute destination path reject;
  \item page-side POST reject;
  \item allowed host + IPv4 és IPv6 működik;
  \item mixed DNS answer egyetlen tiltott címmel fail closed;
  \item forbidden port, TLS identity mismatch és peer-IP mismatch fail closed;
  \item request/response byte limit és timeout enforced;
  \item unsupported upstream transfer/content encoding reject;
  \item cumulative external-I/O exhaustion resource-limit failure;
  \item secret nem folyhat generic outbound JSON-ba.
\end{itemize}

\section{A fejezet lényege}

A külső HTTPS nem utility function. Az RWLang csak \textbf{named, typed, bounded capabilityt} ad; a hálózati destination authority trusted deployment policy marad. Így az SSRF-védelem, effect visibility, TLS policy és resource safety a secure path része.
